% ── Sección: El Prompt ────────────────────────────────────────────────────────
\section{El Prompt}

% ── El Prompt ─────────────────────────────────────────────────────────────────
\begin{frame}[fragile]{El prompt que genero esta presentacion}

  \begin{lstlisting}[basicstyle=\ttfamily\small, breaklines=true]
"example sera una presentacion con el
titulo 'Esta presentacion es toda
vibecodeada', la presentacion nos
explicara la instalacion, uso de pptex
(plug and play latex) y al final el
prompt con el que se hizo la
presentacion"
  \end{lstlisting}

  \vspace{0.6em}
  \begin{exampleblock}{Herramienta}
    \textbf{Claude Code} (Sonnet 4.6) — escrito en la terminal,
    sin tocar un solo archivo \texttt{.tex} manualmente.
  \end{exampleblock}
\end{frame}


% ── Cierre ────────────────────────────────────────────────────────────────────
{
\setbeamercolor{background canvas}{bg=pptexBlue}
\begin{frame}[plain]
  \centering
  \vspace{2.5em}
  {\LARGE\color{white}\textbf{plug \& play LaTeX}}\\[1em]
  {\normalsize\color{white!80}
    \texttt{git clone} \quad \texttt{./pptex build} \quad \texttt{./pptex new slides mi-charla}
  }\\[2em]
  {\small\color{white!60}
    Hecho enteramente con Claude Code (Sonnet 4.6)\\[0.3em]
    \url{github.com/Wilmar3752/pptex}
  }
\end{frame}
}
